\section{The Writer: Confessions, Trinity, the Two Cities, and the Self-Review}

\subsection{The \work{Confessions}}

About 397--400, in the first years of his episcopate, Augustine wrote the
book that created its own genre: thirteen books of prayer in which the
narrative of his life to Monica's death (books 1--9) opens into analyses of
memory, time, and creation (books 10--13). It is autobiography only in
form; in intention it is theology---the demonstration, worked in the
material of one life, that ``you have made us for yourself, and our heart
is restless until it rests in you.'' Its retrospective character is not a
modern discovery but its own repeated confession: it narrates the past ``not
in these words, yet to this effect,'' as remembered by a bishop who reads
God's providence backward through his sins. Its best-remembered sentences
---``Too late did I love You, O Fairness, so ancient, and yet so new! Too
late did I love You! For behold, You were within, and I without, and there
did I seek You'' (\work{Confessions} 10.27.38)---belong to the analytic
books, not the story. Augustine himself, decades on, measured its success
drily: ``which of my smaller works has been able to be more generally and
more agreeably known than the books of my \work{Confessions}?'' (\work{On
the Gift of Perseverance} 53).

\subsection{\work{On the Trinity}}

The \work{De Trinitate}, fifteen books on ``the central core of the
Christian faith,'' occupied him across two decades: the first twelve books
between about 399 and 412, when an unfinished copy escaped and circulated
without his consent, and the whole revised and completed about 420
(Benedict XVI, General Audience, 20 February 2008). Its second half built
the ``psychological'' analogies---memory, understanding, will; the mind
remembering, knowing, and loving God---that governed Latin Trinitarian
theology thereafter. Augustine's own sense of its audience survives in the
letter, quoted in the same catechesis, explaining why he had suspended
dictation of it: ``because they are too demanding and I think that few can
understand them; it is therefore urgent to have more texts which we hope
will be useful to many'' (\work{Letter} 169.1.1, as quoted by Benedict
XVI). The pastor set the systematician's agenda, not the reverse.

\subsection{\work{The City of God}}

In 410 Alaric's Goths entered and sacked Rome, and the shock reached Africa
with the refugees. Pagans said the abandoned gods had abandoned Rome;
shaken Christians half believed them. Augustine's answer, ``an impressive
work crucial to the development of Western political thought and the
Christian theology of history\ldots\ was written between 413 and 426 in 22
books'' at the urging of the same Marcellinus who had judged the Donatists
(Benedict XVI, General Audience, 20 February 2008). Its opening line
announces the scale: ``The glorious city of God is my theme in this work,
which you, my dearest son Marcellinus, suggested, and which is due to you by
my promise'' (\work{City of God} 1, preface). The argument's hinge is the
doctrine of the two loves: ``Accordingly, two cities have been formed by two
loves: the earthly by the love of self, even to the contempt of God; the
heavenly by the love of God, even to the contempt of self'' (\work{City of
God} 14.28). Around that hinge Augustine reordered the moral vocabulary of
antiquity---``it is a brief but true definition of virtue to say, it is the
order of love'' (\work{City of God} 15.22)---and built, in book 19, the
analysis of peace that political theory has quarried ever since: ``The peace
of all things is the tranquillity of order. Order is the distribution which
allots things equal and unequal, each to its own place'' (\work{City of God}
19.13). The work ends where the restless heart was always headed: ``There we
shall rest and see, see and love, love and praise. This is what shall be in
the end without end'' (\work{City of God} 22.30). Quotations here follow the
Dods translation of 1871 and were checked against the exact text artifacts
registered in this repository's source library.

\subsection{The \work{Retractationes}}

Near the end, ``shortly before the time of his death,'' Possidius reports,
``he revised the books which he had dictated and edited''---those written
as layman, presbyter, and bishop---``and whatsoever he found not agreeing
with the ecclesiastical rule, this he himself censured and corrected. Thus
he wrote two volumes whose title is On the Revision of Books''
(\work{Vita} 28). The result, the two books of
\work{Retractationes} of about 427, reviews his works in order,
title by title, marking what he would now say otherwise. Official reception
has singled out the gesture: Benedict XVI called it ``a unique and very
precious literary document but also a teaching of sincerity and
intellectual humility'' (General Audience, 20 February 2008). No other
ancient author left the Church an authorized catalogue of his own
second thoughts; joined to Possidius's \work{Indiculus}, it is the reason
Augustine's corpus survives with its edges defined.
